\section{Explicit fields and arbitrary prefix flags}\label{sec:fields}

For $K\geq2$, let $A$ be the largest power of three strictly smaller
than $K$, so $A<K\leq3A$, and put
\begin{equation}\label{eq:fields}
 \begin{gathered}
 b=4^K,\qquad
 B=\F_2[Z]/\Phi_{9A}(Z),\quad w=6A,\qquad
 \\
 E=\F_2[Y]/\Phi_{27A}(Y),\quad D=18A,
 \qquad Z=Y^3.
 \end{gathered}
\end{equation}
Here and below $\Phi$ is a cyclotomic polynomial, reduced modulo two.
The fields are given in their native monomial bases.

\begin{lemma}\label{lem:fields}
If $m$ is a power of three, then
$\Phi_{3m}(T)=T^{2m}+T^m+1$ is irreducible over $\F_2$.
The chart in \eqref{eq:fields} is an explicit degree-three extension
$E/B$. With $\eta=Y^{9A}$, its relative Frobenius sends
$Y$ to $\eta Y$, where $\eta^2+\eta+1=0$.
\end{lemma}
\begin{proof}
The order of $2$ modulo $3^{j+1}$ is $2\cdot3^j$.
For example, $v_3(4^{3^j}-1)=j+1$, proved inductively by expanding
$(1+3^{j+1}c)^3$, proves the order assertion. Thus the order equals
$\varphi(3m)=2m$, which proves the cyclotomic irreducibility.
Substitution gives $\Phi_{27A}(Y)=\Phi_{9A}(Y^3)$.
The identity $4^{3A}\equiv1+9A\pmod{27A}$ gives the displayed
relative Frobenius. In particular $1,\eta,\eta^2$ are the
specified four-element subfield, with $\eta^2=1+\eta$.
\end{proof}

Write $Q=4^a$ for a power of three $a$, and let $u$ be the native
root of order $9a$ in $\F_{Q^3}$. Then
\[
 \F_{Q^3}=\F_Q\oplus u\F_Q\oplus u^2\F_Q,
 \qquad u^Q=\eta u.
\]
Starting from the basis $(1)$ over $\F_4$, recursively order each
cubic extension basis in three consecutive blocks:
\[
 (\psi_i)_i,\qquad (u\psi_i)_i,\qquad (u^2\psi_i)_i.
\]
Explicitly, let $\mathcal E_1=(0)$ and define ordered exponent lists by
\[
 \mathcal E_{3a}=(3\mathcal E_a,\ 3\mathcal E_a+1,\ 3\mathcal E_a+2).
\]
The lower native root is $u^3$, so the first block $u^{3e}$ is
exactly the previously defined lower-field basis, followed by its
$u$ and $u^2$ multiples. These are ternary digit-reversed native
powers. In the fixed ambient chart the node root is explicitly
$u=Y^{3A/a}$, so $u^{3a}=\eta$ at every node. The recursive order
is therefore consistent under the actual native field inclusions.
Let $W_k$ be
the initial $k$-dimensional $\F_4$ subspace in the smallest such field
containing it. In particular $W_K\subseteq B$ and $1\in W_K$.

For a node $k>1$, take the largest power of three $a<k$ and define
\begin{equation}\label{eq:split}
 (m,r,l)=
 \begin{cases}
 (a,k-a,1),& k\leq2a,\\
 (2a,k-2a,2),& 2a<k\leq3a.
 \end{cases}
 \qquad M=4^m,\quad \ell=4^r.
\end{equation}
Thus $1\leq r\leq a$ and $m\geq k/2$. No prefix is padded to the
next full field dimension.

\begin{lemma}[Sparse low-prefix annihilators]\label{lem:annihilator}
For $Q=4^a$, the annihilator of $W_m$ and the upper image scale are
\[
 \begin{array}{ll}
 P_1(X)=X^Q+X,&c_1=P_1(u)=\eta^2u,\\[2pt]
 P_2(X)=X^{Q^2}+\eta^2X^Q+\eta X,
       &c_2=P_2(u^2)=\eta u^2.
 \end{array}
\]
For the split \eqref{eq:split},
$W_k=W_m\oplus u^lW_r$ and $P_l(x+u^lv)=c_lv$ for
$x\in W_m$, $v\in W_r$.
\end{lemma}
\begin{proof}
$W_a=\F_Q$ and $W_{2a}=\F_Q\oplus u\F_Q$ by the recursive
basis order. The displayed polynomials vanish on these spaces: use
$u^Q=\eta u$ and $u^{Q^2}=\eta^2u$. They are monic and have the
correct degree and nonzero derivative, so these are all their roots.
Every $v\in W_r$ belongs to $\F_Q$, and $Q$-powering fixes $v$.
Consequently $P_l(u^lv)=P_l(u^l)v=c_lv$. Linearity proves the rest.
\end{proof}

The coefficients transformed recursively are arbitrary elements of
the same ambient $E$, of binary width $D$, throughout. The flag
describes evaluation points, not a promise that data lie in a smaller
coefficient field. The bounded case $k=1$ is the four-point
$\F_4$ Vandermonde map, with fixed scalar matrices on native words.
Its constants are $0,1,\eta,\eta^2$; native root shifts and XORs
implement each in $O(D)$ time. This base is uniform in arbitrarily
large ambient width $D$.

